\subsection{ANTICOMMUTATIVE ALGEBRAS AND ALTERNATING ALGEBRAS}

DEFINITION 7. A graded A-algebra E of type Z is called anticommutative if for all non-zero homogeneous elements x, y of E

\[
x y = (- 1) ^ {\deg (x) \deg (y)} y x.\tag{36}
\]

The algebra \(\mathbf{E}\) is called alternating if it is anticommutative and also \(x^{2} = 0\) for every homogeneous element \(x\in \mathbf{E}\) of odd degree.

Remarks. (1) Let \(\mathbf{E}^{+}\) be the graded subalgebra of \(\mathbf{E}\) the direct sum of the \(\mathrm{E}_{2n}\) ( \(n \in \mathbf{Z}\) ); it follows from Definition 7 that if \(\mathbf{E}\) is anticommutative, \(\mathbf{E}^{+}\) is a subalgebra contained in the centre of \(\mathbf{E}\) (and hence commutative).

\begin{enumerate}
\def\labelenumi{(\arabic{enumi})}
\setcounter{enumi}{1}
\item
  Suppose that 2 is not a divisor of 0 in E; then if E is anticommutative E is alternating, since for \(x \in E\) homogeneous and of odd degree, \(x^{2} = -x^{2}\) by (36), whence \(2x^{2} = 0\) and \(x^{2} = 0\) by virtue of the hypothesis.
\item
  We shall study in detail in § 7 important examples of alternating algebras.
\end{enumerate}

Lemma 4. Let E be a graded algebra of type Z and S a set of homogeneous elements \(\neq0\) ; the set F of elements of E all of whose homogeneous components \(x\neq0\) satisfy relation (36) for all \(y\in S\) is a graded subalgebra of E.

It suffices to note that: (1) if \(x'\) , \(x''\) are two homogeneous elements of the same degree \(p\) , \(y\) a homogeneous element of degree \(q\) and \(x'y = (-1)^{pq}yx'\) , \(x''y = (-1)^{pq}yx''\) , then also \((x' + x'')y = (-1)^{pq}y(x' + x'')\) ; (2) if \(x'\) , \(x''\) are two homogeneous elements of respective degrees \(p'\) , \(p'', y\) a homogeneous element of degree \(q\) and \(x'y = (-1)^{p'q}yx'\) , \(x''y = (-1)^{p''q}yx''\) , then

\[
(x ^ {\prime} x ^ {\prime \prime}) y = (- 1) ^ {(p ^ {\prime} + p ^ {\prime \prime}) q} y (x ^ {\prime} x ^ {\prime \prime}).
\]

PROPOSITION 13. Let E be a graded A-algebra of type Z and S a generating system of the algebra E consisting of homogeneous elements \(\neq0\) ; for E to be anticommutative (resp. alternating), it is necessary and sufficient that (36) hold for all \(x\in S\) and \(y\in S\) (resp. that this condition holds and further that \(x^{2}=0\) for all x homogeneous of odd degree belonging to S).

We consider first the case of anticommutative algebras. By Lemma 4, the subalgebra F consisting of the elements all of whose homogeneous components \(x \neq 0\) satisfy (36) for all \(y \in S\) , contains all the elements of S and hence F = E. If now \(F'\) is similarly the subalgebra of E consisting of the elements all of whose homogeneous components \(x \neq 0\) satisfy (36) for every homogeneous element \(y \neq 0\) , it follows from the above that \(F'\) contains all the elements of S and hence \(F' = E\) , which completes the proof of the proposition in this case.

To prove the proposition in the case of alternating algebras, it can be assumed that E is already anticommutative; it is then immediate that every homogeneous element of odd degree in E is of the form \(\sum_{i} z_{i} x_{i}\) , where \(z_{i} \in E^{+}\) and \(x_{i} \in S\) is of odd degree (using the fact that \(E^{+}\) is contained in the centre of E); it follows that \(\left(\sum_{i} z_{i} x_{i}\right)^{2} = \sum_{i} z_{i}^{2} x_{i}^{2} + \sum_{i < j} z_{i} z_{j} (x_{i} x_{j} + x_{j} x_{i}) = 0\) since \(x_{i}^{2} = 0\) by hypothesis and \(x_{i} x_{j} + x_{j} x_{i} = 0\) by (36).

PROPOSITION 14. Let E and F be two graded A-algebras of type Z, both anticommutative (resp. alternating). Then the skew tensor product \(\mathbf{E}^{\mathbf{g}}\otimes_{\mathbf{A}}\mathbf{F}\) (no. 7) is an anticommutative (resp. alternating) algebra.

A generating system of \(E^{g} \otimes_{A} F\) consists of the \(x \otimes y\) , where x (resp. y) is a homogeneous element \(\neq 0\) of E (resp. F). Consider two such elements \(x \otimes y\) , \(x' \otimes y'\) , with \(\deg(x) = p\) , \(\deg(y) = q\) , \(\deg(x') = p'\) , \(\deg(y') = q'\) , so that \(x \otimes y\) is of degree \(p + q\) and \(x' \otimes y'\) of degree \(p' + q'\) . Then by definition (no. 7, formula (27)) and by virtue of (36),

\[
(x \otimes y) (x ^ {\prime} \otimes y ^ {\prime}) = (- 1) ^ {q p ^ {\prime}} (x x ^ {\prime}) \otimes (y y ^ {\prime})
\]

\[
\begin{array}{r l} (x ^ {\prime} \otimes y ^ {\prime}) (x \otimes y) & = (- 1) ^ {p q ^ {\prime}} (x ^ {\prime} x) \otimes (y ^ {\prime} y) \\ & = (- 1) ^ {p q ^ {\prime} + p p ^ {\prime} + q q ^ {\prime}} (x x ^ {\prime}) \otimes (y y ^ {\prime}) \end{array}
\]

and the criterion of Proposition 13 shows that \(\mathbf{E}^{\mathbf{g}}\otimes_{\mathbf{A}}\mathbf{F}\) is anticommutative since \(pq' + p p' + qq' - q p' \equiv (p + q)(p' + q')\) (mod. 2). If further E and F are alternating and \(p + q\) is odd, one of the numbers \(p, q\) is necessarily odd, hence \((x \otimes y)^2 = \pm (x^2) \otimes (y^2) = 0\) and Proposition 13 shows that \(E^g \otimes_A F\) is alternating.

COROLLARY. Let E be an anticommutative (resp. alternating) graded A-algebra of type Z. Then for every ring homomorphism \(\rho: \mathrm{A} \to \mathrm{B}\) , the graded B-algebra \(\rho^{*}(\mathrm{E})\) (no. 8, Remark 2) is anticommutative (resp. alternating).

The ring B with the trivial graduation can be considered as an alternating A-algebra and \(\rho^{*}(\mathbf{E}) = \mathbf{E}^{\mathfrak{g}}\otimes_{\mathbf{A}}\mathbf{B}\) , hence Proposition 14 can be applied.

Remark. Let E be an anticommutative graded A-algebra of type Z. Then the A-linear mapping of \(E \otimes_{A} E\) into E defined by multiplication of E (§1, no. 3) is a homomorphism of the graded A-algebra \(E^{g} \otimes_{A} E\) into E, for in the notation of Proposition 14, in the algebra E,

\[
(x y) (x ^ {\prime} y ^ {\prime}) = (- 1) ^ {q p ^ {\prime}} (x x ^ {\prime}) (y y ^ {\prime}).
\]

